\section{Introduction}

Long-horizon agent systems increasingly combine planning, tool use, memory, verification, and multiple specialized agents. Public benchmarks have improved the field's ability to measure task completion, but production deployments face a broader reliability problem: the runtime must bound cost, detect contract violations, avoid silent failures, contain faulty artifacts, recover without unbounded retry loops, and preserve traces that support diagnosis.

\system reframes an existing multi-agent engineering prototype as a contractual runtime for reliable high-throughput multi-agent workflows. The prior advertising generation workflow remains an example domain plugin. The paper contribution is a reusable runtime abstraction and a stress evaluation harness that can wrap public task environments without changing their task data.

This work studies three research questions. First, can executable contracts over schemas, semantics, dependencies, budgets, and recovery policies turn malformed or off-policy agent outputs into explicit runtime events rather than implicit downstream corruption? Second, can bounded recovery avoid the retry amplification of retry-only execution while preserving a non-zero recovery path under injected production faults? Third, can event-sourced traces make runtime failures auditable enough to support fault localization and reproducible post-hoc diagnosis?

The main matrix is run against the live Qwen3-8B endpoint served by vLLM 0.23.0 and the public \mbox{$\tau^3$}-bench and AgentChangeBench task fixtures. The evidence does not support a simple ``highest task-success'' claim: schema-only validation slightly outperforms the full runtime on raw \mbox{$\tau^3$}-bench task success (15.3\% vs.\ 14.8\%). Instead, the current result supports a narrower runtime claim: \system{} is the only evaluated configuration with non-zero bounded recovery on the main matrix, while retry-only amplifies calls without producing runtime-level diagnosis. Across 1{,}392 live-LLM rows, all four configurations report zero measured silent failures; we therefore treat silent-failure reduction as a metric definition and audit signal, not as standalone evidence that contracts are sufficient.

The key distinction from a benchmark paper is that the workload is not new. tau2-bench / \mbox{$\tau^3$}-bench and AgentChangeBench provide public task environments and fixture metadata~\cite{tau2bench,agentchangebench}. \system contributes the execution substrate around those tasks: contract enforcement, fault localization, bounded recovery, dead-letter handling, and trace-derived reliability metrics.

This paper makes the following contributions:

\begin{itemize}
    \item A contractual execution graph for long-horizon multi-agent workflows, including contracts over artifacts, dependencies, resources, and recovery behavior.
    \item Executable contracts and bounded recovery semantics implemented in a runtime kernel with a scheduler, monitor, artifact store, fault localizer, recovery controller, and event-sourced trace.
    \item A production-stress harness that reuses public benchmark fixtures while running comparable runtime strategies under matched models, tools, fault distributions, concurrency, and recovery budgets.
    \item A trace-derived reliability metric suite that complements fixture-backed task success with dead-letter rate, recovery success, retry amplification, explicit contract violations, and auditability signals.
\end{itemize}

The current artifact includes deterministic smoke tests for development and an audited external main matrix for tau2-bench / \mbox{$\tau^3$}-bench and AgentChangeBench. The paper claims in Section~\ref{sec:experiments} are generated from the external matrix rather than smoke rows.
